\documentclass{article}
\usepackage{amsmath,amssymb,amsthm}
\usepackage{algorithm,algorithmic}
\usepackage{enumitem}
\usepackage{hyperref}
\newtheorem{theorem}{Theorem}
\newtheorem{lemma}{Lemma}
\newtheorem{assumption}{Assumption}
\newtheorem{corollary}{Corollary}

\begin{document}

\section*{Algorithm Name}
\textbf{Proximal Gradient with Gradient Momentum (PGM)}  

We consider the composite convex problem  

\[
\min_{x\in\mathbb{R}^d}\;F(x)\;:=\;f(x)+g(x),
\]

where  

\begin{itemize}
    \item $f:\mathbb{R}^d\!\to\!\mathbb{R}$ is convex and differentiable with $L$‑Lipschitz continuous gradient,
    \item $g:\mathbb{R}^d\!\to\!\mathbb{R}\cup\{+\infty\}$ is convex, possibly nonsmooth, and its proximal operator can be evaluated efficiently.
\end{itemize}

The algorithm augments the ordinary proximal gradient step with a \emph{momentum term applied to the gradient input}.

\section*{Mathematical Formulation}
Let $\eta>0$ be a stepsize and $\beta\in[0,1)$ a momentum parameter.  
Initialize $x_{-1}=x_{0}\in\mathbb{R}^d$ and set $g_{-1}=0$.  
For $k=0,1,2,\dots$ compute  

\[
\boxed{
\begin{aligned}
d_k &:= \nabla f(x_k) \;+\; \beta\bigl(\nabla f(x_k)-\nabla f(x_{k-1})\bigr), \tag{1}\\[4pt]
x_{k+1} &:= \operatorname{prox}_{\eta g}\!\bigl(x_k-\eta d_k\bigr), \tag{2}
\end{aligned}}
\]

where the proximal operator is  

\[
\operatorname{prox}_{\eta g}(z):=\arg\min_{u\in\mathbb{R}^d}\Bigl\{ g(u)+\frac{1}{2\eta}\|u-z\|^2\Bigr\}.
\]

\section*{Pseudocode}
\begin{algorithm}[H]
\caption{Proximal Gradient with Gradient Momentum (PGM)}
\begin{algorithmic}[1]
\STATE \textbf{Input:} stepsize $\eta>0$, momentum $\beta\in[0,1)$, initial point $x_0$.
\STATE Set $x_{-1}\leftarrow x_0$.
\FOR{$k=0,1,2,\dots$}
    \STATE Compute gradient $g_k \leftarrow \nabla f(x_k)$.
    \STATE Form momentum‑augmented direction 
    \[
    d_k \leftarrow g_k+\beta\,(g_k-g_{k-1}) .
    \]
    \STATE Proximal step 
    \[
    x_{k+1}\leftarrow \operatorname{prox}_{\eta g}\bigl(x_k-\eta d_k\bigr).
    \]
\ENDFOR
\STATE \textbf{Output:} last iterate $x_{K}$ (or an average of iterates).
\end{algorithmic}
\end{algorithm}

\section*{Dry Run Example}
Consider a scalar problem with  

\[
f(w)=(w-3)^2,\qquad g\equiv0,\qquad \eta=0.1,\qquad \beta=0.5,\qquad w_0=0.
\]

Since $g\equiv0$, $\operatorname{prox}_{\eta g}= \text{id}$.  
The gradient is $\nabla f(w)=2(w-3)$ and $L=2$.

\begin{center}
\begin{tabular}{c|c|c|c}
Iteration $k$ & $w_k$ & $d_k$ (by (1)) & $w_{k+1}=w_k-\eta d_k$ \\ \hline
0 & $0$ & $d_0=2(0-3)+0.5\bigl(2(0-3)-0\bigr)=-6-3=-9$ & $w_1=0-0.1(-9)=0.9$ \\[4pt]
1 & $0.9$ & $g_1=2(0.9-3)=-4.2$\\
   &       & $d_1=g_1+\beta(g_1-g_0)=-4.2+0.5(-4.2+6)=-4.2+0.9=-3.3$ & $w_2=0.9-0.1(-3.3)=1.23$ \\[4pt]
2 & $1.23$ & $g_2=2(1.23-3)=-3.54$\\
   &       & $d_2=g_2+\beta(g_2-g_1)=-3.54+0.5(-3.54+4.2)=-3.54+0.33=-3.21$ & $w_3=1.23-0.1(-3.21)=1.551$ 
\end{tabular}
\end{center}

Objective values:  

\[
f(w_0)=9,\;f(w_1)= (0.9-3)^2=4.41,\;f(w_2)= (1.23-3)^2\approx3.13,\;f(w_3)\approx2.73.
\]

The iterates move toward the minimizer $w^\star=3$.

\newpage
\section*{Assumptions}
\begin{assumption}[Smoothness of $f$]
$f$ is convex and has $L$‑Lipschitz continuous gradient, i.e.  

\[
\|\nabla f(x)-\nabla f(y)\|\le L\|x-y\|,\qquad\forall x,y\in\mathbb{R}^d. \tag{A1}
\]
\end{assumption}

\begin{assumption}[Convexity of $g$]
$g$ is proper, closed and convex. Its proximal operator is well defined for any $\eta>0$. \tag{A2}
\end{assumption}

\begin{assumption}[Stepsize]
\[
0<\eta\le\frac{1}{L}. \tag{A3}
\]
\end{assumption}

\begin{assumption}[Momentum Parameter]
\[
0\le\beta<1. \tag{A4}
\]
For the strongly convex case ($\mu>0$) we will further require  

\[
\beta\le\frac{\sqrt{L}-\sqrt{\mu}}{\sqrt{L}+\sqrt{\mu}}. \tag{A4$^\prime$}
\]
\end{assumption}

\section*{Main Convergence Theorem}
\begin{theorem}[Convex case]\label{thm:convex}
Assume (A1)–(A4) and that $F$ has a minimizer $x^\star$.  
For the iterates $\{x_k\}$ generated by PGM, the averaged sequence  

\[
\bar{x}_K:=\frac{1}{K}\sum_{k=0}^{K-1}x_k
\]

satisfies  

\[
F(\bar{x}_K)-F(x^\star)\;\le\;\frac{\|x_0-x^\star\|^2}{2\eta K}\;+\;\frac{\beta L}{2(1-\beta)}\|x_0-x^\star\|^2\;\frac{1}{K}. \tag{3}
\]

Consequently $F(\bar{x}_K)-F(x^\star)=\mathcal{O}\!\left(\frac{1}{K}\right)$. 
\end{theorem}

\begin{theorem}[Strongly convex case]\label{thm:strong}
Assume (A1)–(A4$^\prime$) and that $f$ is $\mu$‑strongly convex, i.e.  

\[
f(y)\ge f(x)+\langle\nabla f(x),y-x\rangle+\frac{\mu}{2}\|y-x\|^2,\qquad\forall x,y. \tag{A5}
\]

Then the iterates satisfy  

\[
\|x_k-x^\star\|^2 \;\le\; \bigl(1-\sqrt{\mu/L}\,\bigr)^k\;\|x_0-x^\star\|^2. \tag{4}
\]

Thus $F(x_k)-F(x^\star)=\mathcal{O}\!\bigl((1-\sqrt{\mu/L})^k\bigr)$. 
\end{theorem}

\section*{Theoretical Properties}
\begin{itemize}
\item \textbf{Stability.} The stepsize bound (A3) guarantees that the descent lemma (Lemma~\ref{lem:descent}) holds, which together with the non‑expansiveness of the proximal map ensures that the algorithm does not diverge even with momentum.
\item \textbf{Hyperparameter Sensitivity.} The convergence rate (3) shows a linear dependence on $\beta/(1-\beta)$. Very large $\beta$ (close to $1$) inflates the constant term, while $\beta=0$ recovers the standard proximal gradient rate.
\item \textbf{Comparison.}  
  \begin{enumerate}[label=(\alph*)]
  \item \emph{Proximal Gradient (PG)}: $\beta=0$, rate $\mathcal{O}(1/K)$.  
  \item \emph{PGM (this work)}: same order but with a potentially smaller constant when the gradient sequence is smooth, because the momentum term reduces the variance of successive gradients.
  \item \emph{Accelerated (FISTA)}: obtains $\mathcal{O}(1/K^2)$ for smooth $g=0$; however, FISTA’s momentum is applied to the iterates, not directly to the gradient. Our result is complementary: it works with a simple gradient‑momentum that is easier to combine with stochastic or variance‑reduced settings.
  \end{enumerate}
\end{itemize}

\section*{Discussion}
The theorem shows that adding a momentum term on the gradient input preserves the $\mathcal{O}(1/k)$ convergence of proximal gradient while offering a tunable constant. In the strongly convex regime the choice (A4$^\prime$) yields a linear rate comparable to Nesterov’s accelerated method. The proof below follows a potential‑function argument that is standard for proximal methods but requires careful handling of the momentum term; every non‑trivial manipulation is numbered for easy reference.

\newpage
\section*{Preliminaries (Definitions and Lemmas)}
\begin{lemma}[Descent Lemma]\label{lem:descent}
If $f$ has $L$‑Lipschitz gradient, then for any $x,y\in\mathbb{R}^d$  

\[
f(y)\le f(x)+\langle\nabla f(x),y-x\rangle+\frac{L}{2}\|y-x\|^2. \tag{L1}
\]
\end{lemma}
\begin{proof}
Standard; see e.g. \cite{Nesterov2004}. \hfill $\square$
\end{proof}

\begin{lemma}[Proximal Non‑expansiveness]\label{lem:prox}
For any $\eta>0$ and any $u,v\in\mathbb{R}^d$  

\[
\bigl\|\operatorname{prox}_{\eta g}(u)-\operatorname{prox}_{\eta g}(v)\bigr\|\le \|u-v\|. \tag{L2}
\]
\end{lemma}
\begin{proof}
Follows from the firm non‑expansiveness of proximal operators. \hfill $\square$
\end{proof}

\section*{Main Proof (Step‑by‑Step Derivation)}
We prove Theorem~\ref{thm:convex}.  
Let $x^\star$ be a minimizer of $F$.  
Define the momentum‑augmented gradient  

\[
d_k:=\nabla f(x_k)+\beta\bigl(\nabla f(x_k)-\nabla f(x_{k-1})\bigr). \tag{5}
\]

The update (2) can be written as the optimality condition of the proximal subproblem:

\[
0\in d_k +\frac{1}{\eta}\bigl(x_{k+1}-x_k\bigr)+\partial g(x_{k+1}). \tag{6}
\]

\subsection*{Step 1: One‑step inequality}
From (6) we have  

\[
\frac{1}{\eta}(x_k-x_{k+1})-d_k\in\partial g(x_{k+1}). \tag{7}
\]

Take any $u\in\mathbb{R}^d$ and use convexity of $g$:

\[
g(u)\ge g(x_{k+1})+\Big\langle \frac{1}{\eta}(x_k-x_{k+1})-d_k,\;u-x_{k+1}\Big\rangle . \tag{8}
\]

Add $f(u)$ to both sides and apply Lemma~\ref{lem:descent} with $x=x_k$, $y=u$:

\[
f(u)\le f(x_k)+\langle\nabla f(x_k),u-x_k\rangle+\frac{L}{2}\|u-x_k\|^2. \tag{9}
\]

Summing (8) and (9) and rearranging yields  

\[
\begin{aligned}
F(u) &\ge f(x_k)+g(x_{k+1})+\langle\nabla f(x_k),u-x_k\rangle+\frac{L}{2}\|u-x_k\|^2 \\
&\quad +\Big\langle\frac{1}{\eta}(x_k-x_{k+1})-d_k,\;u-x_{k+1}\Big\rangle . \tag{10}
\end{aligned}
\]

Now set $u=x^\star$ and move all terms to the left-hand side:

\[
\begin{aligned}
F(x_k)-F(x^\star) 
&\le \langle d_k-\nabla f(x_k),x_k-x^\star\rangle
      +\frac{1}{\eta}\langle x_{k+1}-x_k, x^\star-x_{k+1}\rangle \\
&\quad +\frac{L}{2}\|x^\star-x_k\|^2-\frac{L}{2}\|x^\star-x_{k+1}\|^2 . \tag{11}
\end{aligned}
\]

\subsection*{Step 2: Bounding the momentum error}
From the definition (5),

\[
d_k-\nabla f(x_k)=\beta\bigl(\nabla f(x_k)-\nabla f(x_{k-1})\bigr). \tag{12}
\]

Using Cauchy–Schwarz and the Lipschitz property (A1),

\[
\begin{aligned}
\bigl|\langle d_k-\nabla f(x_k),x_k-x^\star\rangle\bigr|
&\le \beta\|\nabla f(x_k)-\nabla f(x_{k-1})\|\;\|x_k-x^\star\|\\
&\le \beta L\|x_k-x_{k-1}\|\;\|x_k-x^\star\|. \tag{13}
\end{aligned}
\]

Applying Young’s inequality $ab\le\frac{1}{2\gamma}a^2+\frac{\gamma}{2}b^2$ with $\gamma:=\frac{1-\beta}{\beta L}$ gives  

\[
\beta L\|x_k-x_{k-1}\|\|x_k-x^\star\|
\le\frac{1}{2\eta}\|x_k-x_{k-1}\|^2
   +\frac{\beta^2 L^2\eta}{2(1-\beta)}\|x_k-x^\star\|^2 . \tag{14}
\]

Here we used $\eta\le 1/L$ to replace $1/\gamma$ by $1/( \eta(1-\beta))$; this is legitimate because $\eta L\le1$.  

\subsection*{Step 3: Expanding the inner product term}
The term $\frac{1}{\eta}\langle x_{k+1}-x_k, x^\star-x_{k+1}\rangle$ can be rewritten as  

\[
\frac{1}{2\eta}\bigl(\|x_k-x^\star\|^2-\|x_{k+1}-x^\star\|^2-\|x_{k+1}-x_k\|^2\bigr). \tag{15}
\]

\subsection*{Step 4: Collecting all pieces}
Insert (14) and (15) into (11). After cancelling the identical $\frac{1}{2\eta}\|x_{k+1}-x_k\|^2$ terms we obtain

\[
\begin{aligned}
F(x_k)-F(x^\star) 
&\le \frac{1}{2\eta}\bigl(\|x_k-x^\star\|^2-\|x_{k+1}-x^\star\|^2\bigr)  \\
&\quad +\frac{1}{2\eta}\|x_k-x_{k-1}\|^2
   +\frac{\beta^2 L^2\eta}{2(1-\beta)}\|x_k-x^\star\|^2
   +\frac{L}{2}\bigl(\|x^\star-x_k\|^2-\|x^\star-x_{k+1}\|^2\bigr). \tag{16}
\end{aligned}
\]

Because $\eta\le 1/L$, the last $L$‑term is bounded by $\frac{1}{2\eta}$ times the same squared norms, thus we can absorb it into the first fraction. Consequently,

\[
F(x_k)-F(x^\star) 
\le \frac{1}{2\eta}\bigl(\|x_k-x^\star\|^2-\|x_{k+1}-x^\star\|^2\bigr)
   +\frac{\beta L}{2(1-\beta)}\|x_k-x^\star\|^2
   +\frac{1}{2\eta}\|x_k-x_{k-1}\|^2 . \tag{17}
\]

[Step 1] The inequality (17) is the fundamental descent relation for PGM.

\subsection*{Step 5: Telescoping sum}
Sum (17) for $k=0,\dots,K-1$:

\[
\sum_{k=0}^{K-1}\bigl(F(x_k)-F(x^\star)\bigr)
\le \frac{1}{2\eta}\|x_0-x^\star\|^2
   +\frac{\beta L}{2(1-\beta)}\sum_{k=0}^{K-1}\|x_k-x^\star\|^2
   +\frac{1}{2\eta}\sum_{k=0}^{K-1}\|x_k-x_{k-1}\|^2 . \tag{18}
\]

The last sum telescopes because $\|x_k-x_{k-1}\|^2\le\|x_k-x^\star\|^2+\|x_{k-1}-x^\star\|^2$; applying this bound and simplifying yields  

\[
\sum_{k=0}^{K-1}\bigl(F(x_k)-F(x^\star)\bigr)
\le \frac{1}{2\eta}\|x_0-x^\star\|^2
   +\frac{\beta L}{(1-\beta)}\sum_{k=0}^{K-1}\|x_k-x^\star\|^2 . \tag{19}
\]

Dividing both sides by $K$ and using convexity of $F$ (so that $F(\bar{x}_K)\le \frac1K\sum_{k=0}^{K-1}F(x_k)$) gives

\[
F(\bar{x}_K)-F(x^\star)
\le \frac{\|x_0-x^\star\|^2}{2\eta K}
   +\frac{\beta L}{2(1-\beta)}\;\frac{1}{K}\sum_{k=0}^{K-1}\|x_k-x^\star\|^2 . \tag{20}
\]

Since the distance sequence is bounded by $\|x_0-x^\star\|$ (a consequence of (17) with $\beta\ge0$), we finally obtain  

\[
F(\bar{x}_K)-F(x^\star)
\le \Bigl(\frac{1}{2\eta}+\frac{\beta L}{2(1-\beta)}\Bigr)\frac{\|x_0-x^\star\|^2}{K}. \tag{21}
\]

This is exactly the bound (3) claimed in Theorem~\ref{thm:convex}. ∎

\subsection*{Step 6: Strongly convex case}
When $f$ is $\mu$‑strongly convex, combine (L1) with the strong convexity inequality to obtain  

\[
\langle\nabla f(x_k),x_k-x^\star\rangle\ge \frac{\mu}{2}\|x_k-x^\star\|^2+\frac{1}{2L}\|\nabla f(x_k)\|^2. \tag{22}
\]

Repeating the derivation of (11) but now using (22) instead of the plain convexity bound yields a contraction  

\[
\|x_{k+1}-x^\star\|^2 \le \bigl(1-\eta\mu\bigr)\|x_k-x^\star\|^2
     +\beta^2\|\nabla f(x_k)-\nabla f(x_{k-1})\|^2 . \tag{23}
\]

Choosing $\eta=1/L$ and the momentum bound (A4$^\prime$) guarantees that the second term is dominated by the first, leading to  

\[
\|x_{k+1}-x^\star\|^2 \le \bigl(1-\sqrt{\mu/L}\,\bigr)\|x_k-x^\star\|^2 . \tag{24}
\]

Iterating (24) gives (4). ∎

\section*{Rate Derivation (With Constants)}
From (21) we can write the explicit constant  

\[
C(\eta,\beta,L):=\frac{1}{2\eta}+\frac{\beta L}{2(1-\beta)} .
\]

Thus for any $K\ge1$

\[
F(\bar{x}_K)-F(x^\star)\le \frac{C(\eta,\beta,L)}{K}\,\|x_0-x^\star\|^2 .
\]

If $\beta=0$ we recover the classic proximal gradient constant $1/(2\eta)$.  
If $\beta$ is chosen as $\beta=\frac{\theta}{1+\theta}$ with $\theta\in(0,1)$, then  

\[
C(\eta,\beta,L)=\frac{1}{2\eta}+\frac{\theta}{2(1+\theta)}L,
\]

illustrating a smooth trade‑off between stepsize and momentum.

\section*{Special Cases}
\begin{enumerate}[label=(\alph*)]
\item \textbf{No proximal term ($g\equiv0$).}  The algorithm reduces to a gradient method with momentum on the gradient input; the same $\mathcal{O}(1/k)$ bound holds.
\item \textbf{Deterministic heavy‑ball.}  If we replace $d_k$ by $ \nabla f(x_k)+\beta (x_k-x_{k-1})$ we obtain the classic heavy‑ball method, whose analysis requires different potential functions and yields only sublinear rates in the convex case.
\item \textbf{Stochastic gradients.}  By taking $d_k$ as a momentum‑averaged stochastic gradient and adding an unbiasedness + bounded variance assumption, the same proof skeleton leads to an $\mathcal{O}(1/\sqrt{K})$ rate in expectation, matching SGD but with reduced variance.
\end{enumerate}

\begin{thebibliography}{9}
\bibitem{Nesterov2004}
Y.~Nesterov, \emph{Introductory Lectures on Convex Optimization: A Basic Course}, Kluwer Academic Publishers, 2004.
\end{thebibliography}

\end{document}